\section{Compiling classical formulas with seeded controls}
\label{sec:classical-synthesis}

The preceding section supplies a bounded classical formula whose rooted
energy is close to the SK optimum.  We now implement any formula in that
class by quantum controls with independent Gaussian longitudinal seeds.
A \emph{seeded circuit} is a finite word in the cost, the uniform mixer,
and the gates
\[
 \exp\!\left(-\ii t\sum_i s_iZ_i\right),\qquad t\in\R,
\]
where the same independent standard Gaussian family \((s_i)\) is used
throughout the word and averaged in its value.  Here
\(s_i=s_{\mathrm{raw},i}\) is the raw seed of
\eqref{eq:raw-seed-reencoding}: every colored target formula is evaluated
at \(C(s_i),S(s_i)\).  Both the color and the scalar Gaussian input
are therefore functions of the seed driving this single longitudinal
control.  All angles are deterministic.
Section~\ref{sec:seed-simulation} will remove this extra control after
the present construction has selected one finite word.

We use the masks, aggregations \(A_{RS}\), and seed-only exclusion
formulas of Definition~\ref{def:exclusion-formulas}.  In particular,
every transmitted nonlinear function and every output angle is bounded,
although aggregate fields need not be bounded.  A mask may recur in
different branches of an expression tree, but not along a single active
branch.  We retain this distinction throughout the proof: formulas that
share descendants are not assumed independent.

All approximations in this section take place in the infinite-degree
rooted representation, inside a fixed surrounding circuit.  We first
construct elementary masked controls, then use unused helper spins to
implement one aggregation, and finally compile a finite expression tree
from its outermost operations inward.  The helper contraction is the
main step.  Its proof controls both the visible aggregation and the
changes it induces inside every lower formula.

\subsection{The synthesis statement}

Fix a finite collection of such formulas \(\theta_R\), odd and
seed-only, one for each output mask \(R\).  Their ideal operation is
\begin{equation}
 U_{\mathrm{cl}}=\prod_i \exp(\ii\theta_iY_i),
 \qquad m_i=\sin(2\theta_i).
 \label{eq:ideal-classical-rotations}
\end{equation}
The product is locally defined: an output at a vertex uses only a
fixed-radius neighborhood.  The factors commute because the angles
are seed-only.  Reserve disjoint helper masks, absent from all these
formulas, and set the output angle to zero on them.  One helper mask
per aggregation depth is sufficient.  Let \(H\) denote their union
and let \(p_H\) be its probability.

\begin{theorem}[Classical synthesis]
\label{thm:classical-synthesis}
For every finite collection of bounded smooth odd seed-only exclusion
formulas of Definition~\ref{def:exclusion-formulas}, with disjoint
positive-probability helper masks, one for each positive aggregation
depth, absent from all their dependencies and with output angle zero on
their union \(H\), and every \(\eta>0\), there is a finite seeded
circuit whose insertion vector \(v\) satisfies
\begin{equation}
 \norm{\1_{H^c}(v-v_{\mathrm{cl}})}<\eta,
 \label{eq:synthesis-insertion}
\end{equation}
where \(v_{\mathrm{cl}}\) is the insertion vector of
\eqref{eq:ideal-classical-rotations}.  In particular its value obeys
\begin{equation}
 \abs{\cE(v)-\cE(m)}\le 2\eta+2\sqrt{p_H}.
 \label{eq:synthesis-energy}
\end{equation}
The circuit and all its angles are chosen independently of the degree,
system size, and disorder.  Helper probabilities can be arbitrarily
small but are fixed and positive when the construction is performed.
\end{theorem}

The proof will establish convergence in a fixed surrounding circuit,
not only convergence on a freshly initialized input.  This is needed
because a formula is called repeatedly inside commutator products.

We use the smooth spaces of Section~\ref{sec:foundations}.  At the
root let
\[
 W=1+\dd\Gamma(\Lambda),\qquad
 \norm{T}_{a\leftarrow b}=\norm{W^aTW^{-b}},\qquad
 \norm{d}_{a}=\norm{\Lambda^a d}.
\]
The bounded-seed version of the weights is used until the last
compilation step.  Thus the weights do not differentiate or multiply
the root seed.  In particular they commute with root masks on the odd
subspace.  A family of operators is \emph{uniformly smooth} if for
each \(a\) there are \(b,C\), independent of the member of the
family, with \(\norm{T}_{a\leftarrow b}\le C\); the same convention
applies to adjoints and to parameter derivatives.  For a parameter
\(y\), a polynomial factor in \(1+|y|\) is allowed.  A statement
\(T=O_{\mathrm{sm}}(\varepsilon)\) means that these bounds hold
with \(C\varepsilon\) for every fixed output power and every fixed
jet.  Conditional versions mean the essential supremum in the root
seed.  The constants may depend on the fixed formulas, time intervals,
and positive helper probabilities.

The conditional bounds in Theorem~\ref{thm:gate-calculus} imply the
following useful rule.  If \(T\vac\) has bounded conditional smooth
norms and \(R\) is a root mask, then
\begin{equation}
 \norm{W^a\1_R T\vac}\le C_a\sqrt{p_R}.
 \label{eq:synthesis-masked-column-bound}
\end{equation}
For odd columns, the removal of the vacuum in the definition of
\(\Lambda\) does not change this estimate or the orthogonality of
columns supported on disjoint root masks.

\subsection{Elementary masked controls and matrix products}
\label{sec:elementary-masked-controls}

We first identify the elementary gates from which the helper
operation will be built.  These gates need only the original Gaussian
longitudinal control, parity-preserving local rotations, and the cost.

We first treat the following gates as exact:
\begin{enumerate}
\item local \(X\) rotations with bounded even seed coefficients,
      and local \(Y\) or \(Z\) rotations with bounded odd seed
      coefficients;
\item interactions between two disjoint masks in any pair of odd axes
      in the \(YZ\) plane.
\end{enumerate}
We justify at the end of this subsection that they lie in the closure
of finite seeded circuits.  Working with exact elementary masked gates
in the interim keeps unused helpers exactly pinned.

We will repeatedly use pointwise matrix sums and commutators.  Here
is the topology of that use.  If \(F(x),G(x)\) are Hermitian
two-by-two matrices with smooth polynomially bounded derivatives,
the usual product and group-commutator formulas converge on compact
sets with every fixed number of derivatives.  Their derivative bounds
are polynomial in \(x\), uniformly in the subdivision.  To verify the
last, write the group-commutator loop as
\[
 V_h(x)=e^{\ii hF(x)}e^{\ii hG(x)}
                e^{-\ii hF(x)}e^{-\ii hG(x)}.
\]
Twice integrating its mixed derivative in the two loop times shows
\[
 \norm{\partial_x^\alpha(V_h-1)}
     \le C_\alpha h^2(1+|x|)^{r_\alpha},\qquad |h|\le1.
\]
For a differentiated product of \(N\) such loops, each differentiated
factor contributes \(O(h^2)\); at a fixed derivative order the sum
is a fixed polynomial in \(Nh^2\).  Undifferentiated factors are
unitary.  Thus \(h=N^{-1/2}\) gives the claimed uniform bounds.
The Taylor error of a loop is \(O(h^3)\), with all fixed jets on a
compact set, so telescoping \(N\) factors proves convergence there.
The sum formula is identical with \(Nh\) in place of \(Nh^2\).
Theorem~\ref{thm:gate-calculus} therefore permits these limits inside
any fixed surrounding circuit.

When such a matrix formula is used simultaneously at many vertices,
the coefficient algebra must be common and commuting, and exclude
the pivot mask.  Under that condition the entire finite-tree gate is
a product of exactly those two-by-two formulas.  Its reinterpretation
at a new pivot is permitted only after an exact finite-tree identity
has been established.  These identities are also valid on every cut
tree, so the rooted maps agree.

To obtain the local elementary controls from the raw seeded controls,
write \(s=s_{\mathrm{raw}}\) and conjugate \(X\) by
\(e^{\ii t sZ}\).  Every local coefficient in the colored formulas,
including every mask, is a measurable function of this same \(s\)
by \eqref{eq:raw-seed-reencoding}.  Taking sums and differences
gives \(\cos(2ts)X\) and \(\sin(2ts)Y\); a constant \(X\)
rotation changes \(Y\) to \(Z\).  Finite even and odd trigonometric
sums are thus available, together with their exponentials.  Bounded
measurable even or odd seed functions are limits in Gaussian measure
of uniformly bounded continuous functions with the same parity.
After taking a subsequence, convergence is almost everywhere.  A box
cutoff, periodic extension, and Fejer approximation give trigonometric
approximants with the same uniform bound and parity.  Local seed
derivatives are not part of the smooth topology, so bounded dominated
convergence and Theorem~\ref{thm:gate-calculus} apply in a fixed
surrounding circuit.  If intermediate expressions use the raw linear
seed, its polynomial moments are controlled by the seed-weighted
version of that theorem.

We isolate interactions by sign averaging, a standard refocusing
device for Ising couplings; see Jones and
Knill~\cite[Section~II]{JonesKnill1999}.  Refine the masks to a finite
partition \(R_1,\ldots,R_M\).  For a sign vector
\(\epsilon\in\{-1,1\}^M\), masked \(X\) flips conjugate the
ordinary cost to
\[
 C_\epsilon=D^{-1/2}\sum_{\{i,j\}}
       \epsilon_{R(i)}\epsilon_{R(j)}Z_iZ_j.
\]
For distinct \(a,b\), sign orthogonality gives the finite-tree
identity
\[
 C_{R_aR_b}=2^{-M}\sum_\epsilon\epsilon_a\epsilon_b C_\epsilon.
\]
All the terms on the right commute, so their exponential implements
the desired subcost exactly when the masked flips are treated as
exact.  Separate masked \(X\) rotations put its two endpoints in any
specified axes in the \(YZ\) plane.  This proves the claimed
elementary compilation without introducing interactions within a
mask.

\subsection{A contraction through unused spins}

The elementary local rotations constructed so far depend only on the
seed at the rotated vertex.  The next construction extends them to a
neighbor sum of a fixed lower formula.  It uses a disjoint set of
spins kept in their \(X=1\) state, and its error tends to zero when the sender
mask is divided into sufficiently small parts.

Let \(I,J,K\) be disjoint masks, with \(K\) an unused helper of
probability \(p_K>0\).  Let \(f_j\), on \(J\), be a bounded odd
seed-only exclusion formula whose cross-edge dependencies avoid
\(I,J,K\).  Other reserved helpers are avoided as well.  Temporarily
allow exact rotations by \(f\) on \(J\), including their restrictions
to further local seed masks.  These ideal calls will remain in the
finite approximants constructed here.  At the end of the section we
will remove them by compiling the full formula from its outermost
operations inward.  We shall obtain
\begin{equation}
 H_{I,J,f}=\sum_{i\in I}P_i(A_{IJ}f)_i,
 \qquad P\in\{Y,Z\}.
 \label{eq:desired-propagation}
\end{equation}

Partition \(J\) into finitely many even seed masks \(J_l\), with
probabilities \(p_l\).  The formula \(f\) is kept fixed when this
partition is refined.  The commutators below sum paths of the form
\begin{equation}
 \begin{gathered}
 \begin{array}{ccccccc}
 i\in I&-&j\in J_l&-&k\in K&-&j'\in J_l\\
 &&&&&&\downarrow\\[-2pt]
 &&&&&&f_{j'}
 \end{array}\\
 j'\ne j:\ \text{simple path},\qquad
 j'=j:\ i-j-k-j\ \text{(backtrack)}.
 \end{gathered}
 \label{eq:helper-path-schematic}
\end{equation}
The downward arrow denotes the lower calculation of \(f\), whose
cross-edge dependencies avoid all three displayed masks.  In the
backtrack, \(f\) is attached at \(j'=j\); the two sender spin
factors cancel and the pinned helper supplies its polarization.
The simple paths contribute an error that
disappears from nonhelper insertions as the partition is refined.
We must check their effects at every root cut, including cuts
inside the lower formula \(f\).

Write \(A_{ij}=D^{-1/2}\1_{i\sim j}\) and
\(b_j=\sum_{i\in I}A_{ji}P_i\).  Pointwise commutators on \(J_l\)
give
\begin{align}
 H_l&=\sum_{i\in I,j\in J_l,k\in K}
       A_{ij}A_{jk}P_iX_jZ_k,\\
 G_l&=\sum_{j\in J_l,k\in K}A_{jk}X_jf_jY_k.
 \label{eq:two-helper-generators}
\end{align}
For example, \((2\ii)^{-1}[\sum A_{ij}P_iY_j,
\sum A_{jk}Z_jZ_k]=H_l\), and
\((2\ii)^{-1}[\sum f_jY_j,\sum A_{jk}Z_jY_k]=G_l\).
These are exact identities: the only noncommuting factors meet on
\(J_l\), and the seed-only coefficients commute with every spin.

At the next pivot, \(K\), the two coefficients are
\[
 B_{k,l}=\sum_{j\in J_l}A_{kj}X_jb_j,
 \qquad F_{k,l}=\sum_{j\in J_l}A_{kj}X_jf_j.
\]
They commute with each other and with all spins on \(K\).  Their
calculation paths avoid \(K\), as required by the gate calculus.
Consequently \(-(2\ii)^{-1}[H_l,G_l]
=\sum_{k\in K}X_k B_{k,l}F_{k,l}\).  Summing and rescaling gives
\begin{equation}
 \widetilde H_{\mathcal P}
  =\frac1{p_K}\sum_l\sum_{
       i\in I,\,j,j'\in J_l,\,k\in K}
       P_i A_{ij}X_j A_{jk}X_k A_{kj'}X_{j'}f_{j'}.
 \label{eq:helper-contraction-generator}
\end{equation}
All its summands commute.  At every fixed partition
\(\mathcal P\), its exponential is therefore in the closure of
finite products of elementary masked gates and calls to \(f\).
This statement holds on the full spin space and in every fixed
surrounding circuit; it has not yet used helper pinning.

We now specify the space on which pinning is used.  On a finite
light cone with vertex set \(V\), define the orthogonal projection
\begin{equation}
 \Pi_{K,V}=\prod_{v\in V}
       \left[1-\1_K(s_v)\frac{1-X_v}{2}\right].
 \label{eq:helper-pinning-projection}
\end{equation}
Its range fixes every helper-labeled spin in \(\ket+\), at the root
and throughout its descendants, while retaining all seed coordinates.
For the degree limit, let
\[
 \begin{gathered}
 \mathsf k_K=\operatorname{Ran}\!\left(1-\1_K(s)\frac{1-X}{2}\right),\\
 \cH_{K,0}=\mathsf k_K,\qquad
 \cH_{K,h+1}=\mathsf k_K\otimes\Gamma_s(\cH_{K,h}^{\circ}),
 \end{gathered}
\]
where the centered subspace is taken relative to the same vacuum
\(\vac_h\) as before.  These spaces are subspaces of \(\cH_h\).
The occupation and empty-descendant embeddings restrict to them,
as do the bounded-seed weights, so their smooth norms are inherited
from the original recursion.  We call this the \emph{compressed space}.
Helper-labeled descendant spaces remain present; only their helper
spins are pinned, and the seed law is unchanged.

The ideal propagations, the exact generators
\eqref{eq:helper-contraction-generator}, and surrounding gates avoiding
\(K\) preserve this range.  The observable \(\1_{K^c}(s_o)Z_o\)
at the root \(o\) also preserves it.  Thus nonhelper insertions may
be computed in the compressed space.  This restriction is different
from multiplication of an insertion vector by \(\1_{K^c}\), which
selects only the root label and imposes no condition on descendant
spins.  The individual commutator factors need not preserve pinning;
their finite products will be chosen after the exact contraction.

\begin{proposition}[Simultaneous helper contraction]
\label{prop:helper-contraction}
Consider a fixed finite sequence in the parity-preserving,
root-seed-commuting gate-calculus class, containing finitely many
propagations \eqref{eq:desired-propagation}, possibly with different
\(I,J,f,P\), and otherwise fixed gates.  All its gates and classical
parameters avoid a common helper \(K\).  Start with its spins pinned
at \(X=1\).  Replace every indicated propagation by
\eqref{eq:helper-contraction-generator}, and refine their sender
partitions so that their largest part probability tends to zero.
Then all insertion vectors off \(K\) converge to those of the
original sequence.  The root maps and their adjoints have uniform
conditional smooth bounds on the compressed space, and converge
strongly there up to scalar phases depending only on the root label.
The same estimates hold on cut subtrees.  Responses through the fixed
lower formulas for \(f\) have the uniform jet estimates proved below;
no convergence of arbitrary boundary jets in the growing list of
partition-dependent top channels is asserted.
\end{proposition}

The contraction must control every root cut, including cuts inside
\(f\), uniformly as the sender partition is refined.  The next two
estimates provide the analytic bounds for its proof.  Most cuts use a
fixed number of fields.  At a \(K\)-root, however, each part
\(J_l\) supplies a pair of incoming fields, and the proof will
subtract their same-neighbor covariance.  The second estimate controls
this growing sum by the orthogonality of the masks.

\subsection{Estimates and proof of the contraction}

\begin{lemma}[A fixed number of commuting fields]
\label{lem:fixed-field-exponentials}
Let \(Q_j=q(d_j)\), \(1\le j\le k\), commute, where \(k\) is
fixed and \(\norm{d_j}_a\le C_a\) for every \(a\).  Let
\(F_s:\R^k\to\R\) have polynomial bounds on every derivative,
uniformly in the local seed \(s\).  For bounded \(t\), the operators
\(\exp(\ii tF_s(Q))\), their adjoints, and their time derivatives
are uniformly smooth.  Auxiliary parameter jets obey the same conclusion
when the weighted column jets and the joint derivatives of \(F\)
have the corresponding uniform polynomial bounds.  The constants
require no lower bound on the
Gram matrix of the \(d_j\)'s.  The same assertion holds with a fixed
root control axis, by conditioning on its two eigenvalues.
\end{lemma}

\begin{proof}
The Weyl estimates in Lemma~\ref{lem:weyl-bounds}, applied to
\(\sum_j\xi_jd_j\), give, for every \(a\),
\[
 \norm{W^a e^{\ii\xi\cdot Q}\zeta}
 \le C_a(1+|\xi|)^{r_a}\norm{W^{b_a}\zeta}.
\]
Only upper bounds for weighted column norms enter this inequality.
For an integer \(N\), put
\(g_{s,t}(x)=e^{\ii tF_s(x)}(1+|x|^2)^{-N}\).
Fix the desired output power.  First choose enough derivatives that
integration by parts gives an integrable Fourier bound with weight
\((1+|\xi|)^{r_a}\), and then choose \(N\) large enough that those
derivatives are integrable, uniformly in \(s,t\).  The polynomial
derivative assumptions permit this order of choices.  Functional
calculus gives, on smooth vectors,
\[
 e^{\ii tF_s(Q)}
  =\int_{\R^k}\widehat g_{s,t}(\xi)
        e^{\ii\xi\cdot Q}\,\dd\xi\,(1+|Q|^2)^N,
\]
with the chosen Fourier normalization absorbed in \(\widehat g\).
The polynomial on the right has a fixed smooth power bound by the
creation and annihilation estimates.  Fourier integration therefore
proves the assertion.  Differentiating in any additional parameters
also differentiates the fields and the polynomial factor.  If
\(v(y)=\sum_j\xi_jd_j(y)\), the central commutator relation gives
\begin{equation}
 \partial_y e^{\ii q(v(y))}
   =e^{\ii q(v(y))}
       \left(\ii q(\partial_yv)
                 +\ii\im\braket{v}{\partial_yv}\right).
 \label{eq:moving-column-weyl-derivative}
\end{equation}
Repeated derivatives are the same Weyl operator times finite ordered
products of derivative fields and scalar Gram pairings.  Each field
insertion has a smooth bound polynomial in \(|\xi|\), and every
scalar pairing is quadratic in \(\xi\), by the weighted column-jet
hypotheses.  The Leibniz rule gives the same finite-power bounds for
derivatives of \((1+|Q|^2)^N\); its differentiated fields need not
commute with the original list.  Increasing the Fourier decay and
input weight powers just chosen proves the jet assertions.  No
change of field coordinates, and hence no inverse Gram matrix,
occurs.
\end{proof}

\begin{lemma}[Quadratics with many orthogonal columns]
\label{lem:orthogonal-quadratics}
Let \(R_l\) be finitely many pairwise disjoint root masks, with
\(p_l=\Pr(R_l)\). Suppose \(d_l,e_l\) are supported on \(R_l\)
and are real-commuting, meaning that all fields \(q(d_l),q(e_l)\)
commute pairwise. Suppose also that for every \(a\)
\[
 \norm{d_l}_a+\norm{e_l}_a\le C_a\sqrt{p_l}.
\]
Set
\[
 Q=\sum_l\left(q(d_l)q(e_l)-\braket{d_l}{e_l}\right),
 \qquad \rho=\left(\sum_l p_l^2\right)^{1/2}.
\]
Then \(Q=O_{\mathrm{sm}}(\rho)\).  Its exponentials and adjoints
are uniformly smooth for bounded times, independently of the number
of columns.  They converge to the identity strongly on smooth vectors
as \(\rho\to0\), and in every fixed lower smooth norm on families
with all higher smooth norms bounded.
\end{lemma}

\begin{proof}
Normal ordering gives a pair-creation kernel, its adjoint, and a
one-particle kernel:
\[
 Q=a^*(\textstyle\sum_l d_l\otimes e_l)
    +a(\textstyle\sum_l e_l\otimes d_l)
    +\dd\Gamma\!\left(\textstyle\sum_l
          |d_l\rangle\langle e_l|+|e_l\rangle\langle d_l|\right),
\]
where the first kernel is understood symmetrically.  Orthogonality of
different masks gives
\[
 \left\|\sum_l\Lambda^a d_l\otimes\Lambda^b e_l\right\|_{\rm HS}
 +\left\|\sum_l
       |\Lambda^a d_l\rangle\langle\Lambda^b e_l|\right\|_{\rm HS}
 \le C_{a,b}\rho.
\]
The same inequalities hold with \(d,e\) interchanged.  Pair creation
by a Hilbert--Schmidt kernel has norm at most
\(\sqrt{(n+1)(n+2)}\) times its kernel norm on the \(n\)-particle
sector; the corresponding estimates for annihilation and a
one-particle kernel are bounded by a constant times \(n+1\).
Commuting powers of \(W\) through these operators inserts powers of
\(\Lambda\) on their kernel indices.  This proves the asserted
smooth smallness.

For completeness, this smallness of the generator alone is not used
to assert an exponential power bound.  Every iterated commutator
\(\ad_W^j(Q)\) is another quadratic with uniformly bounded
Hilbert--Schmidt kernels of the displayed form.  Consequently
\[
 \norm{\ad_W^j(Q)\zeta}\le C_j\norm{W\zeta}.
\]
Writing \([W^a,Q]\) with these commutators on the left yields
\(\norm{[W^a,Q]\zeta}\le C_a\norm{W^a\zeta}\) for positive
integer \(a\).  On smooth vectors, differentiation of
\(\norm{W^ae^{\ii tQ}\zeta}^2\) and Gronwall's inequality now give
a bound independent of the column count.  At a fixed finite column
count these differentiations are justified by
Lemma~\ref{lem:fixed-field-exponentials}; its constants are not used
in Gronwall's bound.  Equivalently one may first truncate the weights
and pass to the limit using these same inequalities.

Finally, spectral calculus gives
\[
 \norm{(e^{\ii tQ}-1)\zeta}\le |t|\norm{Q\zeta}
       \le C|t|\rho\norm{W^b\zeta}.
\]
The uniform higher-power bounds and interpolation prove the last
assertion.  The same conclusions hold for \(U^*QU\) when \(U\) is
unitary and both \(U\) and \(U^*\) have uniform smooth bounds.  They
are uniform for families of unitaries with common such bounds: indeed,
\(e^{\ii tU^*QU}=U^*e^{\ii tQ}U\), and the smooth power losses
compose.
\end{proof}

\begin{proof}[Proof of Proposition~\ref{prop:helper-contraction}]
We prove a single replacement with arbitrary preceding root maps
from the gate calculus: they preserve the joint parity, commute with
root seed multiplication, and have uniform conditional smooth bounds.
These hypotheses are preserved by all the gates used here.  For the
sender partition put
\[
 \rho=\left(\sum_l p_l^2\right)^{1/2},\qquad
 \varepsilon=\max_l\sqrt{p_l}.
\]
We will estimate the Hamiltonian correction at each root cut, prove
uniform smooth bounds for its exponential, and then iterate the
replacement.  At helper roots a scalar phase remains; we will check
that it cancels from the child signals before the next replacement.
The partition changes only the terminal coefficients of the lower
calculation: the formula \(f\) and its finite channel list stay fixed.

On the range of \eqref{eq:helper-pinning-projection}, and hence in its
recursive occupation realization, we may set \(X_k=1\) at every
\(k\in K\).  The preservation facts just established apply to the
exact replacement and all its surrounding gates.

\paragraph{The relative root operation.}
Fix the partition and introduce an unscaled real time \(t\), keeping
the preceding circuit fixed as \(t\) varies.  Write \(H_D\) for the desired Hamiltonian and \(\widetilde H_D\) for the
complete contraction Hamiltonian.  On a padded rooted tree \(T\),
retain target terms through a depth \(d\) larger than the influence
distance.  On each child component \(T_a\), use the same formula with
zero parent input and target depth \(d-1\).  All parameter formulas
are computed on the padded tree.  For either Hamiltonian \(F_D\), put
\[
 L_{F,D}=F_D(T)-\sum_a F_D(T_a).
\]
All terms in the two Hamiltonians and their branch baselines commute:
they are diagonal in the same physical site axes.  Factor the preceding
circuit as \(U_D=\mathcal U_Du_D\), where
\(\mathcal U_D=\bigotimes_aU_{a,D}\).  The new child circuits are
\(e^{\ii tF_D(T_a)}U_{a,D}\).  Factoring their product from
\(e^{\ii tF_D(T)}U_D\), as in
\eqref{eq:exclusion-root-update}, gives exactly
\[
 u_{F,D}(t)=\mathcal U_D^*e^{\ii tL_{F,D}}\mathcal U_Du_D.
\]
Consequently the relative root maps satisfy
\begin{equation}\label{eq:helper-relative-root-group}
 \mathcal T_{\mathcal P,D}(t)
 :=u_{H,D}(t)^*u_{\widetilde H,D}(t)
 =U_D^*e^{\ii t(L_{\widetilde H,D}-L_{H,D})}U_D.
\end{equation}
This is a unitary group in \(t\), whose exponent is a root cut
difference.  The zero-input branch Hamiltonians have been subtracted;
no limit of the extensive Hamiltonian difference is being asserted.
The identity holds on the compressed space as well.

At this fixed partition, Theorem~\ref{thm:gate-calculus} applies to
both complete operations with \(t\) as an external parameter.  The
complete contraction is the gate at the \(K\) pivot, with coefficient
\(p_K^{-1}\sum_l B_{k,l}F_{k,l}\).  Their local matrices have polynomial bounds for all mixed time and input
derivatives, locally uniformly in \(t\).  The root maps, their adjoints,
and their time derivatives therefore have smooth degree limits.
Passing to the limit in \eqref{eq:helper-relative-root-group} gives a
unitary group \(\mathcal T_{\mathcal P}(t)\): its group law and
inverses follow from the finite-degree identities, and time-derivative
convergence gives strong continuity.  Its self-adjoint generator
\(A_{\mathcal P}\) is defined by this group.  Differentiation at zero
shows that every smooth vector is in its domain and that, under the
occupation identifications,
\begin{equation}\label{eq:helper-total-generator-action}
 A_{\mathcal P}\zeta
 =\lim_{D\to\infty}
 U_D^*(L_{\widetilde H,D}-L_{H,D})U_D\zeta.
\end{equation}

We now compute this derivative by separating diagonal paths, simple
paths, and lower-formula responses.  This is an algebraic decomposition
of the finite-degree cut difference, not an application of the gate
calculus to each separated Hamiltonian.  All contributions to
\eqref{eq:helper-total-generator-action} use the preceding full frame.
Thus top-path expressions first computed in the child frame receive
one final conjugation by \(u\); the lower-response formulas already
display it.  The smooth bounds are preserved by this conjugation.
The diagonal average and its helper-root scalar are included before
any phase is removed.

\paragraph{The diagonal paths.}
The terms \(j=j'\) in
\eqref{eq:helper-contraction-generator} equal
\begin{equation}
 \sum_{i\in I,j\in J}A_{ij}P_i f_j
       \frac{\#\{k\sim j:k\in K\}}{Dp_K}.
 \label{eq:diagonal-helper-paths}
\end{equation}
Away from the incidence lost at a \(K\)-root cut, the last factor
tends to one, with all smooth bounds required in the gate calculus:
it is a normalized average of bounded local seed masks.  This remains
true after conjugation by preceding maps with uniform smooth bounds.
One can see the rate on the symmetric embedding by decomposing that
average into its vacuum scalar, vacuum row and column, and occupied
block.  The centered row and column are \(O(D^{-1/2})\), and the
occupied block is \(O(D^{-1})\) with a loss of one number power.

If the removed root lies in \(K\), the lost incidence contributes
\((Dp_K)^{-1}\sum_{j\sim o}\1_J(j)b_jf_j\).
The average has a scalar vacuum limit.  Thus it contributes only a
phase conditional on the root's membership in \(K\).  Below we also
recover this scalar as the same-neighbor contraction of the full
\(K\)-root expression.  No estimate discards an uncentered
normalized sum.

\paragraph{The nonbacktracking paths at top masks.}
Remove the terms \(j=j'\).  A remaining path
\(i,j,k,j'\) is simple on a tree.  At a terminal \(j'\in J_l\),
looking away from its parent, its extra coefficient is
\begin{equation}
 R_{j',l}=\frac1{p_K}
   \sum_{k\sim j',\,k\in K} A_{j'k}
   \sum_{j\sim k,\,j\in J_l,\,j\ne j'}
         A_{kj}X_jb_j.
 \label{eq:extra-terminal-coefficient}
\end{equation}
Here all inner fields are computed in descendant trees.  Let
\(d_l\) be the child vacuum column of \(\1_{J_l}Xb\), in its
preceding Heisenberg frame.  It is centered, since it is odd, and
\eqref{eq:synthesis-masked-column-bound} gives
\(\norm{d_l}_a\le C_a\sqrt{p_l}\).
The first normalized lift is \(q(d_l)\).  Masking its next root on
\(K\), conjugating by the preceding map, and lifting again gives
\begin{equation}
 R_{j',l}=O_{\mathrm{sm}}(\sqrt{p_l})
 \quad\hbox{conditionally at }j'\in J_l.
 \label{eq:terminal-smallness}
\end{equation}
The constants include \(p_K^{-1}\), which is fixed.  The reversed
coefficient, with \(f\) in place of \(b\), has the same bound.
This reasoning uses centered lifts at both steps; the parity of the
transmitted fields supplies their centering.

At a root in \(I\), the final incoming column is the combination
over \(J_l\) of the small coefficients just estimated.  Its squared
weighted norm is bounded by
\(C\sum_l p_l^2\), since conditioning on \(J_l\) contributes
one factor \(p_l\) and \eqref{eq:terminal-smallness} contributes
another.  The extra root Hamiltonian is its normalized lift times
the root axis \(P\), and is \(O_{\mathrm{sm}}(\rho)\).

At a root in \(J_l\), the two possible occurrences of that root,
as \(j\) or \(j'\), give a bounded smooth factor \(f\) or the
field \(b\) times one of the small coefficients in
\eqref{eq:terminal-smallness}.  The two factors use distinct
first-step masks.  Thus there is no further same-neighbor term to
subtract.  The root correction is
\(O_{\mathrm{sm}}(\sqrt{p_l})\), uniformly conditionally on the
root seed in this part.  It is a smooth function of a fixed number of
commuting fields with uniformly bounded weighted columns.

At a root in \(K\), let \(d_l,e_l\) be the incoming child columns
of \(\1_{J_l}Xb\) and \(\1_{J_l}Xf\).  All their fields commute.
The root contribution of the full top-path sum is
\(p_K^{-1}\sum_l q(d_l)q(e_l)\).  The terms using the same neighbor
are the limit of the average
\[
 \frac1{Dp_K}\sum_{j\sim o}\1_J(j)b_jf_j
       \longrightarrow
       \kappa=\frac1{p_K}\sum_l\braket{d_l}{e_l}.
\]
In particular, this is an average on one branch, not the product of
two independent lifts.  After subtracting it, the correction is
\begin{equation}
 Q_K=\frac1{p_K}\sum_l
       \left(q(d_l)q(e_l)-\braket{d_l}{e_l}\right).
 \label{eq:helper-root-quadratic}
\end{equation}
The weighted columns belonging to different parts are orthogonal.
Lemma~\ref{lem:orthogonal-quadratics} gives
\(Q_K=O_{\mathrm{sm}}(\rho)\) and uniform exponential bounds
independent of the number of parts.  The scalar \(\kappa\) is the
harmless local phase identified from the diagonal paths.

This list accounts for all top-path cuts.  If the root belongs to
\(I,J,K\), a top path not containing it cannot change after the cut:
the dependencies of \(f\) exclude all three masks.  At a \(J\)
root, its own \(f\) uses only descendants.  A root outside both the
top masks and the lower calculation has no correction.  It remains
to check cuts inside the lower computation of \(f\).

\paragraph{Responses through the lower formula.}
We have accounted for cuts through the four-vertex helper paths.  A cut
can also change the value of \(f\) at their terminal vertices.  The
following recursion keeps those changes, including mixed responses
between channels, in the estimate.

First work in the physical frame.  Write
\(H_{\rm off}(T)=\sum_{l,j'\in J_l}\omega_{j',l}f_{j'}^T\),
where \(\omega_{j',l}=X_{j'}R_{j',l}\) initially denotes its full
top-path coefficient.  If a cut root \(v\) has a lower mask, it
cannot lie on any of the top paths.  Each coefficient is consequently
unchanged and lies wholly in one cut component.  For the child
components \(T_a\), the exact finite-tree identity is
\begin{equation}
 H_{\rm off}(T)-\sum_aH_{\rm off}(T_a)
   =\sum_{a,l,j'\in T_a\cap J_l}\omega_{j',l}
          \bigl(f_{j'}^T-f_{j'}^{T_a}\bigr).
 \label{eq:physical-lower-cut-identity}
\end{equation}
Only terminals within the fixed influence radius contribute.  At an
active terminal \(j'\in J_l\), a boundary input into its lower
computation therefore changes the Hamiltonian by
\begin{equation}
 X_{j'}R_{j',l}\bigl(f_{j'}(y)-f_{j'}(0)\bigr).
 \label{eq:terminal-response}
\end{equation}
Whenever this difference is active, its parent lies in a lower mask,
outside \(I,J,K\).  Consequently every first step of
\(R_{j',l}\) goes into a descendant branch.  That coefficient has
no boundary response of its own on this active path.  If the parent
is in a top mask, the boundary input into \(f\) is absent, and
\eqref{eq:terminal-response} is zero.  This observation ensures that
using the downward expression
\eqref{eq:extra-terminal-coefficient} loses no term.

Here is the precise difference recursion.  At a root \(v\) of the
lower calculation let \(m(y)\) be its fixed finite list of outgoing
channel values.  Let \(\mathcal H_v^{[d]}(y)\) be the sum of the
weighted terminal terms with terminal distance at most \(d\),
evaluated with boundary input \(y\), and put
\(\Delta_v^{[d]}(y)=\mathcal H_v^{[d]}(y)-\mathcal H_v^{[d]}(0)\).
The coefficients are computed downward, as just justified; the
terminal formulas themselves use all required sites.  In the
physical frame the recursion is a local term
\eqref{eq:terminal-response}, when the root is terminal, plus
\begin{equation}
 \sum_{a\text{ child}}
 \left[
   \Delta_a^{[d-1]}(D^{-1/2}m(y))
   -\Delta_a^{[d-1]}(D^{-1/2}m(0))
 \right].
 \label{eq:additive-response-recursion}
\end{equation}
Only channels of the fixed \(f\)-calculation are included.  At the
actual output root, where there is no parent input, the root
Hamiltonian correction is obtained from the local term and the
child sum evaluated at \(D^{-1/2}m(0)\), without this subtraction.
The definition counts distance to the terminal vertex; the top
coefficient beyond that vertex has already been computed downward.
Taking \(d\) beyond the finite influence range gives exactly the
Hamiltonian difference across the cut.  Indeed the identical
recursion with \(\mathcal H_a\) first holds by partitioning
terminals among children, and replacing
\(\mathcal H_a(z)=\Delta_a(z)+\mathcal H_a(0)\) cancels the
constant terms.  This also explains the argument \(m(0)\) in the
second summand.

Now factor the preceding circuit as
\(U_v=(\bigotimes_aU_a)u_v\).  These preceding maps are fixed in
the boundary variable of the new Hamiltonian.  Simultaneous
conjugation of the physical identity gives
\begin{equation}
 \widehat\Delta_v(y)=u_v^*\left[T_v(y)+\sum_a
   \left\{\widehat\Delta_a(D^{-1/2}m(y))
             -\widehat\Delta_a(D^{-1/2}m(0))\right\}\right]u_v,
 \label{eq:conjugated-lower-cut-identity}
\end{equation}
where \(\widehat\Delta_a=U_a^*\Delta_aU_a\), and \(T_v(y)\) is
the local terminal term \eqref{eq:terminal-response}, zero at a
nonterminal root. The direct term and messages are interpreted in
the same child frames.  There is no
derivative of \(U_a\) or \(u_v\).  Although a preceding Heisenberg
image of the top coefficient may appear to have larger support, its
boundary independence follows from the physical identity before
this fixed conjugation.  Repeated lower masks in different expression
branches cause no additional terms: these are identities for the
whole nonlinear function of all incoming channels, so their mixed
responses are retained.

The required parity covariance is explicit: if \(\Pi\) is the joint
parity, then \(\Pi\omega\Pi=-\omega\) and
\(\Pi f(y)\Pi=-f(-y)\).  Thus
\(\Pi\{\omega[f(y)-f(0)]\}\Pi=\omega[f(-y)-f(0)]\), and
the recursion preserves this covariance.  We pass the algebraic
recursion \eqref{eq:additive-response-recursion} to the limit by the
centered normalized-sum calculation used in the proof of
Corollary~\ref{cor:additive-response}.  At a fixed partition, the
terminal coefficients and their jets are finite products and smooth
compositions of centered odd normalized sums, together with the
separately estimated bounded-mask averages normalized by \(D^{-1}\).
The normalized-sum bounds and the conditional bounds for preceding
maps give uniform degree bounds in every required higher smooth norm;
these constants may depend on the fixed partition.  They propagate
through the fixed-depth difference recursion.  Taylor expansion in the
boundary inputs therefore retains the centered first response and the
scalar second response, while its third-order remainder carries
\(D\cdot D^{-3/2}\).  Commuting substitution handles the outgoing
message arguments.  These are termwise limits in the derivative of
the complete relative family, not a unitary calculus for the separated
Hamiltonians.  If \(D_\alpha,D_{\alpha\beta}\) are the first two input derivatives
of a child response at zero, set
\[
 r_\alpha=D_\alpha\vac,\qquad
 c_{\alpha\beta}=\bra{\vac}D_{\alpha\beta}\ket{\vac}.
\]
Its contribution is
\begin{align}
 &\sum_\alpha\bigl(m_\alpha(y)-m_\alpha(0)\bigr)q(r_\alpha)
 \nonumber\\[-2pt]
 &\quad+\frac12\sum_{\alpha,\beta}
 \bigl(m_\alpha(y)m_\beta(y)
        -m_\alpha(0)m_\beta(0)\bigr)c_{\alpha\beta},
 \label{eq:limiting-additive-response}
\end{align}
followed by the preceding root conjugation.  At a nonterminal lower
root with no parent input, the Hamiltonian correction instead uses
the unsubtracted child contribution
\[
 u_v^*\left[
 \sum_\alpha m_\alpha(0)q(r_\alpha)
 +\frac12\sum_{\alpha,\beta}
       m_\alpha(0)m_\beta(0)c_{\alpha\beta}
 \right]u_v.
\]
This is the root correction specified after
\eqref{eq:additive-response-recursion}, rather than the value at
\(y=0\) of the boundary difference, which is zero.  Both expressions
use the same response columns and scalar second responses.
First derivatives are odd and hence centered.  The same calculation
applies to every fixed input jet.  The partition estimates below are
made after this degree limit.

At depth zero, \eqref{eq:terminal-smallness} and the derivative
bounds of \(f\) make every fixed response jet
\(O_{\mathrm{sm}}(\varepsilon)\), conditionally in the local
seed.  More explicitly, the sum of terminal jets over \(l\) is a
direct integral on the disjoint events \(J_l\), so its conditional
bound is the maximum of their \(C\sqrt{p_l}\) bounds, not their
sum.  Suppose the induction holds at depth \(d-1\).  All weighted norms of
\(r_\alpha\) and all \(|c_{\alpha\beta}|\) are then bounded by
\(C\varepsilon\).  The functions \(m\) and their derivatives have
fixed polynomial bounds.  Multiplication, differentiation, and the
preceding root conjugation in
\eqref{eq:limiting-additive-response} therefore give the same
smallness at depth \(d\), with a possible finite loss of powers.
This induction has fixed finite depth and a fixed number of channels,
both independent of the number of parts \(J_l\).  It requires no
independence between different lower channels.  In particular it
applies when two expression branches reuse some lower masks.  At a
nonterminal lower root, if there are \(k\) outgoing channels and \(s\) ordinary Gaussian
inputs used to compute \(m\), the root correction uses at most
\(s+k\) field directions.  Each \(r_\alpha\) is a single child
vector containing the whole lower calculation; its Hilbert-space
support need not be finite dimensional.  The scalar second
responses contribute no new directions.

The diagonal degree-factor errors have the same response recursion
with terminal coefficient
\[
 b_{j'}\left(\frac{\#\{k\sim j':k\in K\}}{Dp_K}-1\right)
\]
in place of \(X_{j'}R_{j',l}\).  On an active lower response path
the parent is not in \(K\) or \(I\), so this is again a downward
coefficient.  Its limiting smooth bounds vanish by the average
estimate already proved.  Thus all its lower responses vanish as
well.

\paragraph{Uniform topology and iteration.}
After the degree limit, the diagonal degree-factor errors vanish.
The preceding computations identify the total action in
\eqref{eq:helper-total-generator-action}.  At a helper root its scalar
part is \(\sigma_{\mathcal P}(s)=\kappa\1_K(s)\), with
\[
 |\kappa|\le p_K^{-1}\sum_l\|d_l\|\,\|e_l\|
             \le C p_J/p_K.
\]
This bounded root-seed multiplication commutes with every root map.
Hence
\[
 \mathcal S_{\mathcal P}(t)
 =e^{-\ii t\sigma_{\mathcal P}}\mathcal T_{\mathcal P}(t)
\]
is a strongly continuous unitary group.  Its self-adjoint generator
\(R_{\mathcal P}\) acts on smooth vectors by the total remaining
correction.  The four cut cases, in the common preceding frame, give
\[
 \begin{array}{c|l}
 \text{cut-root mask}&\text{remaining correction}\\ \hline
 I&O_{\mathrm{sm}}(\rho)\quad\text{(combined column)}\\
 J_l&O_{\mathrm{sm}}(\sqrt{p_l})\quad\text{(conditional on }J_l\text{)}\\
 K&u^*Q_Ku=O_{\mathrm{sm}}(\rho)\quad\text{(after its scalar phase)}\\
 \text{lower mask}&O_{\mathrm{sm}}(\varepsilon)
                  \quad\text{(fixed lower channels)}
 \end{array}
\]
Since the parts partition the fixed sender mask,
\[
 \rho^2=\sum_l p_l^2\le p_J\max_l p_l,
 \qquad \varepsilon^2=\max_l p_l.
\]
Both quantities tend to zero under the refinement hypothesis, with
\(f\) and \(p_K>0\) fixed.

To obtain uniform smooth bounds for this actual group, consider the
full field expression for the remaining correction at each root.  At
an \(I\)-root it uses a bounded number of combined columns.  At a
\(J_l\)-root include the fields for both possible sender occurrences.
At a lower root use all ordinary channel fields and response columns
in the unsubtracted root formula.  These current-generator fields
commute: before preceding conjugation they belong to the same physical
seed and fixed-axis algebra.  Thus the total expression is a smooth
function of a fixed number of commuting fields and, when present, the
fixed root axis, with the uniform polynomial derivative bounds already
proved.  At a \(K\)-root it is the complete centered quadratic
\(Q_K\).  Define \(V_{\mathcal P}(t)\) by the joint spectral
calculus of these expressions, as a direct integral over the root seed,
with the final conjugation by \(u\) wherever it is not already
displayed.  Lemma~\ref{lem:fixed-field-exponentials} gives its uniform
conditional smooth bounds in the fixed-field cases;
Lemma~\ref{lem:orthogonal-quadratics} gives the count-uniform bound at
\(K\).  These bounds concern the total expression at each root.

We verify that \(V_{\mathcal P}\) is the actual relative group
\(\mathcal S_{\mathcal P}\).  At fixed partition the field-calculus
group preserves the smooth domain and has smooth time derivatives
there, with
\(\partial_tV_{\mathcal P}(t)\zeta
 =\ii R_{\mathcal P}V_{\mathcal P}(t)\zeta\)
by the action identified in \eqref{eq:helper-total-generator-action}.
For a smooth vector
\(\zeta\), therefore, \(V_{\mathcal P}(s)\zeta\) remains in the
domain of \(R_{\mathcal P}\), and its image under that generator
varies continuously by the smooth bounds.  Differentiating gives
\[
 \frac{d}{ds}\left[
 \mathcal S_{\mathcal P}(t-s)V_{\mathcal P}(s)\zeta
 \right]=0.
\]
Evaluation at \(s=0,t\) proves
\(\mathcal S_{\mathcal P}(t)\zeta=V_{\mathcal P}(t)\zeta\), and
density extends the identity to the Hilbert space.  Thus the uniform
smooth bounds just proved apply to the relative root group itself;
no separate closure of the formal sum of corrections is chosen.

Spectral calculus for its generator now yields
\[
 \|[\mathcal S_{\mathcal P}(t)-1]\zeta\|
 \le |t|\|R_{\mathcal P}\zeta\|\longrightarrow0.
\]
Uniform higher smooth bounds upgrade this to convergence in every
fixed lower graph norm; the same holds for adjoints.  The relative
factor multiplies the desired root map on the right, so
\[
 u_{\widetilde H,\mathcal P}(t)
 =u_H(t)\mathcal T_{\mathcal P}(t)
 =e^{\ii t\sigma_{\mathcal P}}u_H(t)\mathcal S_{\mathcal P}(t).
\]
This proves convergence of the actual root maps up to their root-seed
phase, with uniform conditional smooth bounds.  The same argument
applies on cut subtrees, and the fixed lower response jets retain the
estimates proved above.

A root-seed phase cancels in each descendant Heisenberg conjugation
before its vacuum column is lifted; helper-labeled descendant spaces
need not be removed.  Starting from the fixed initial preceding maps,
apply this single-replacement argument successively through the fixed
finite sequence.  The continuity part of
Theorem~\ref{thm:gate-calculus} applies at each step with the conditional
smooth bounds just established.  It gives insertion convergence off
\(K\) and proves the simultaneous contraction statement.
\end{proof}

\subsection{Pointwise functions and the order of compilation}

Proposition~\ref{prop:helper-contraction} implements an aggregation
whenever its lower formula is still available as an ideal call.  To
remove all ideal calls without losing unused helper spins, we combine
it with pointwise synthesis and process the formula from its outermost
aggregations inward.

Propagation and elementary local controls generate the required
pointwise nonlinearities.  Suppose rotations by odd fields
\(h_1,\ldots,h_k\) are available on the same target mask.  Same-axis
rotations give their real linear combinations.  A constant local
\(X\) rotation interchanges the two odd axes.  For an even local
coefficient \(a(s)\), conjugating \(a(s)X\) by a rotation with odd
angle \(h\), and taking the difference at opposite angles, gives
\(a(s)\sin(2h)Y\), up to a sign convention.  For an odd local
coefficient \(c(s)\), conjugating \(c(s)Y\) by the same \(Z\)
rotation and taking the sum gives \(c(s)\cos(2h)Y\).
Thus every finite trigonometric expression
\begin{equation}
 \sum_\nu a_\nu(s)\sin(\xi_\nu\cdot h)
       +c_\nu(s)\cos(\xi_\nu\cdot h),
 \qquad a_\nu\text{ even},\quad c_\nu\text{ odd},
 \label{eq:parity-trigonometric-synthesis}
\end{equation}
can be implemented as an odd-axis Hamiltonian by pointwise matrix
sums.  This covers joint parity in the local seed and aggregate
arguments; it does not require multiplying arbitrary even fields.

The same construction applies to bounded smooth odd functions of a
finite list of odd aggregate inputs.  To see the required closure
precisely, cut off on larger symmetric boxes, periodically extend,
and approximate the smooth extension by its Fourier polynomials.
The cutoff can be chosen even, and parity projection gives exactly
the coefficient parities in
\eqref{eq:parity-trigonometric-synthesis}.  Uniform bounds for
derivatives in the aggregate variables are preserved on each fixed
compact set.  Choose successive Fourier truncations accurate for
increasing derivative orders.  The cutoffs and periodic extensions
can be chosen with polynomial derivative bounds independent of the
box: inside the box these follow from the original bounds, and
outside its fundamental cell the absolute value of the reduced
coordinate is no larger than that of the original coordinate.
Fourier errors in the finitely many required derivative norms may
be made at most one.  A diagonal choice gives the polynomial
domination and local all-jet convergence used by
Theorem~\ref{thm:gate-calculus}.  Alternatively, if only the terminal
insertion is needed, first truncate the finitely many input laws and
use their smooth moment bounds.  Local seed coefficients remain
bounded measurable functions throughout.

We finish by specifying the nesting of all limits.  This is also
where helper pinning is used.  Define expression rank by
\[
 \operatorname{rank}(g(s))=0,\qquad
 \operatorname{rank}(F(h_1,\ldots,h_k))=\max_j\operatorname{rank}(h_j),
 \qquad
 \operatorname{rank}(A_{RS}f)=1+\operatorname{rank}(f).
\]
Local masking does not change rank.  Reserve a distinct helper
\(K_r\) for each positive rank.  Every original formula avoids
every helper.  Before processing rank \(r\), all unresolved calls
are classical formulas with this property; every already produced
elementary gate acts only on original masks and previously used
helpers \(K_q\), \(q>r\).  Thus the current surrounding word has
no target, spin control, or nonzero emitted seed channel on \(K_r\).
Its spins are exactly pinned, even if rotations on original output
masks have already occurred.  The estimates above allow arbitrary
preceding states in the gate-calculus class on those original masks.

\begin{enumerate}
\item At the largest remaining rank, replace each pointwise
nonlinearity by the pointwise synthesis just described.  Take these
limits successively in the fixed surrounding circuit.  After fixing
their finite products to any desired accuracy, the remaining calls
at rank \(r\) are rotations by explicit propagated odd inputs
\(A_{IJ}f\), where \(\operatorname{rank}(f)\le r-1\).

\item Apply Proposition~\ref{prop:helper-contraction} to all those
propagations simultaneously, using the reserved helper \(K_r\).
It is unused by every other exact gate at this stage and is
therefore pinned.  Choose finite sender subdivisions sufficiently
fine for the desired accuracy.  Sender subdivisions change only
the local calling mask; they do not refine the lower formulas or
their channel lists.

\item With that subdivision fixed, implement all the commutators
and sums defining \eqref{eq:helper-contraction-generator} by finite
pointwise products.  Their convergence holds on the full space, so
these finite products may temporarily rotate the helper away from
\(X=1\).  Pinning was used in the preceding contraction limit,
not during this inner approximation.

\item The resulting finite word contains only lower-rank formula
calls and elementary masked gates.  Repeat the same procedure at
the next rank, using its new helper.  The previously produced
finite word may act on higher helpers, but it avoids this new one
exactly.  Thus the pinning hypothesis holds anew.  Source formulas
retain all ancestor exclusions, including those introduced by their
particular calling masks. A sender subdivision only restricts the
local calling mask: exclusion of the original sender mask implies
exclusion of each submask. The commutators create further occurrences
of the same lower formula, without adding a cross-edge dependency or
introducing a control on any still-unused helper.
\end{enumerate}

This induction terminates because the formula rank is finite.
At each stage there are only finitely many occurrences in the
already fixed surrounding word.  Accuracies can therefore be chosen
successively from the outside inward.  No uniform estimate in the
number of commutator factors, the number of formula occurrences
created by an earlier stage, or inverse helper density is asserted
or needed.  The only uniformity in a growing sender subdivision is
the one explicitly proved in
Proposition~\ref{prop:helper-contraction}.

At termination, replace the elementary masked gates by their finite
seeded-circuit approximations from
Section~\ref{sec:elementary-masked-controls}.  All
mask-controlled operations are compiled at this final stage.
Theorem~\ref{thm:gate-calculus} permits the successive replacements
inside the fixed finite remaining word, including the raw Gaussian
seed control with its moment weight.  Exact helper pinning is no
longer required after the contraction limits have been fixed.
This constructs a finite seeded word with arbitrarily small
insertion error off the helpers.

\begin{proof}[Proof of Theorem~\ref{thm:classical-synthesis}]
The finite nested construction proves
\eqref{eq:synthesis-insertion}.  To identify the target value,
the ideal seed-only rotations give
\[
 v_{\mathrm{cl}}
  =m\ket{+}+\cos(2\theta)\ket{-},
\]
where both coefficients are classical functions of seeds; all
descendant spins remain in their pointed spin vectors.  In
\(\braket{v_{\mathrm{cl}}}{\cL v_{\mathrm{cl}}}\), the second
summand on the left is orthogonal to the range of \(\cL\).  Its
shifted version on the right has a child spin excitation and is
orthogonal to the first, purely classical summand on the left.
Consequently \(\cE(v_{\mathrm{cl}})=\cE(m)\).  The same identity
holds after omitting helper insertions, because \(m=0\) there.

For any circuit preserving root seed labels,
\(\norm{\1_Hv}=\sqrt{p_H}\): the mask commutes with the circuit
and with \(Z\), and \(Z\) is unitary.  Since \(\cL\) is an
isometry, Proposition~\ref{prop:energy} gives
\[
 \abs{\cE(v)-\cE(\1_{H^c}v)}\le2\sqrt{p_H}.
\]
Both projected insertion vectors have norm at most one, so
\eqref{eq:synthesis-insertion} changes their energy by at most
\(2\eta\).  This proves \eqref{eq:synthesis-energy}.  All
approximation limits above compare already-defined fixed-word tree
values.  After choosing finite approximants, the resulting seeded
word and all its parameters are independent of the degree, system
size, and disorder.
\end{proof}
